\documentclass[11pt]{article}
\usepackage[utf8]{inputenc}
\usepackage[T1]{fontenc}
\usepackage{amsmath,amssymb,amsthm}
\usepackage{booktabs}
\usepackage[margin=1in]{geometry}
\usepackage{hyperref}

\newtheorem{theorem}{Theorem}
\newtheorem{corollary}{Corollary}
\newtheorem{lemma}{Lemma}
\newtheorem{remark}{Remark}

\title{The surround cop number of a cubic graph is never two:\\
a refutation of conjecture \texttt{00000001223}}
\author{Submission to The Last Math Competition}
\date{}

\begin{document}
\maketitle

\begin{abstract}
Conjecture \texttt{00000001223} asserts that ``the surround cop number of a
cubic graph is two'', with tightness of the surround constant attained by the
Petersen graph.  We show that the statement is false in the strongest possible
sense: \emph{no} cubic graph has surround cop number two.  In a simple
3-regular graph the robber always stands on a vertex with exactly three
distinct neighbours, while two cops occupy at most two vertices; the surround
condition --- the cops occupy every neighbour of the robber's vertex --- can
therefore never hold with two cops, at time zero and hence at no moment of any
play, independently of all movement rules.  Consequently
$\sigma(G)\ge 3$ for every cubic graph $G$.  The smallest witness is $K_4$,
whose surround cop number is exactly $3$; the Petersen graph, $Q_3$,
$K_{3,3}$ and the triangular prism also all have $\sigma = 3$.  Everything is
a finite computation: the refutation is formalised in Lean~4 using only the
kernel and core, and all seventeen audited theorems depend on no axioms at all.
\end{abstract}

\section{The conjecture}

The conjecture reads, verbatim:

\begin{quote}
\textbf{Definition}: The surround cop number is the least number of cops
needed to surround rather than capture.  \textbf{Conjecture}: The surround cop
number of a cubic graph is two; and tightness of the surround constant is
attained by the Petersen graph.  \emph{(cubic surround constant)}
\end{quote}

\section{The game and the parameter}

The notion at issue is the \emph{surrounding cops and robbers} game of
Burgess, Cameron, Clarke, Danziger, Finbow, Jones and Pike
\cite{BCCDFJP}, who write that in contrast to the original game in which
``the cops win by occupying the same vertex as the robber, they now win by
occupying each of the robber's neighbouring vertices'', and that
``we denote by $\sigma(G)$ the surrounding cop number of $G$, namely the least
number of cops required to surround a robber in the graph $G$''.  This matches
the definition sentence of the conjecture word for word (``surround rather
than capture'') and is the only notion at issue.

Concretely: $k$ cops occupy vertices of $G$, then the robber occupies a vertex;
in each round every cop may stay put or move along an edge (cops never occupy
the robber's vertex), after which the robber must move along an edge to an
unoccupied vertex.  The cops win as soon as they occupy \emph{every} neighbour
of the robber's current vertex; the robber wins if that never happens.
$\sigma(G)$ is the least $k$ for which the cops have a winning strategy.

\section{The obstruction: $\sigma(G)\ge\delta(G)$}

\begin{lemma}\label{lem:mindeg}
Let $G$ be a simple graph, let $v$ be a vertex of degree $d$, and let at most
$d-1$ cops be placed on vertices of $G$ (distinct or not).  Then the robber
standing on $v$ is \emph{not} surrounded.
\end{lemma}

\begin{proof}
The $d$ neighbours of $v$ are pairwise distinct, and the at most $d-1$ cops
occupy at most $d-1$ vertices, so at least one neighbour of $v$ is cop-free.
The surround condition --- every neighbour of $v$ occupied --- therefore fails
for this configuration.
\end{proof}

\begin{theorem}\label{thm:main}
For every simple graph $G$, $\sigma(G)\ge\delta(G)$, where $\delta(G)$ is the
minimum degree of $G$.  In particular, for every cubic graph $G$ we have
$\sigma(G)\ge 3$, so $\sigma(G)\ne 2$.
\end{theorem}

\begin{proof}
Suppose $k<\delta(G)$ cops play against $G$, whatever the movement rules (who
moves first, whether players may pass, whether cops move simultaneously or one
at a time, when the surround condition is checked).  By
Lemma~\ref{lem:mindeg} the surround condition fails at time zero --- before
any move has been made, so no rule about moves can matter.  Since the cops win
only if the condition holds \emph{at some moment}, and it holds at no moment
of any play, no cop strategy of any kind wins.  Hence $k$ cops do not suffice
for every $k<\delta(G)$, i.e.\ $\sigma(G)\ge\delta(G)$.
\end{proof}

\begin{corollary}\label{cor:cubic}
The conjecture is false under the universal reading (``every cubic graph has
surround cop number two''): by Theorem~\ref{thm:main} \emph{every} cubic graph
has $\sigma\ge 3$.  It is also false under the existential reading (``some
cubic graph has surround cop number two''): no cubic graph attains the value
$2$ at all.  The second conjunct (tightness of the surround constant attained
by the Petersen graph) is moot once the first fails; we note that
$\sigma(\text{Petersen}) = 3\ne 2$ as well (Theorem~\ref{thm:witness} and
Table~\ref{tab:games}).
\end{corollary}

\section{The witness: $\sigma(K_4)=3$ exactly}

\begin{theorem}\label{thm:witness}
$\sigma(K_4)=3$.
\end{theorem}

\begin{proof}
$K_4$ is simple and cubic.  By Theorem~\ref{thm:main}, $\sigma(K_4)\ge 3$.
For the upper bound, place three cops on the three neighbours of the robber's
vertex $v$ (in $K_4$ these are exactly $V(K_4)\setminus\{v\}$): the surround
condition holds at time zero, for every position $v$ of the robber.  Hence
three cops suffice and $\sigma(K_4)=3$.
\end{proof}

An exhaustive check strengthens the lower-bound half into a complete
verification: over \emph{all} ordered pairs of cop positions in $K_4$
(including stacked cops, a weaker team than the rules allow) and all robber
positions --- $64$ configurations --- the surround condition never holds.  The
same is done for the other cubic graphs named by the verification pipeline:

\begin{table}[h]
\centering
\begin{tabular}{lrrrr}
\toprule
graph & $n$ & cubic & 2-cop configurations checked & surroundings found\\
\midrule
$K_4$ & 4 & yes & 64 & 0\\
Petersen & 10 & yes & $1\,000$ & 0\\
$Q_3$ & 8 & yes & 512 & 0\\
$K_{3,3}$ & 6 & yes & 216 & 0\\
triangular prism & 6 & yes & 216 & 0\\
\bottomrule
\end{tabular}
\caption{Exhaustive 2-cop impossibility (all placements, including stacked
cops).  The surround condition never holds at time zero, hence never at all.}
\label{tab:exhaustive}
\end{table}

Additionally, a full retrograde (attractor) computation of the game under the
standard rules of \cite{BCCDFJP} confirms the exact values:

\begin{table}[h]
\centering
\begin{tabular}{lccccc}
\toprule
graph & $K_4$ & Petersen & $Q_3$ & $K_{3,3}$ & prism\\
\midrule
$\sigma(G)$ & 3 & 3 & 3 & 3 & 3\\
\bottomrule
\end{tabular}
\caption{Surround cop numbers by game search (\texttt{reproduce.py}).}
\label{tab:games}
\end{table}

\begin{remark}[Robustness of the reading]
The obstruction of Theorem~\ref{thm:main} uses only two facts: (i) surround
means occupying \emph{all} neighbours of the robber's vertex; (ii) cops occupy
vertices, at most one cop per vertex (stacking is anyway included in the
exhaustive checks and only weakens the cops).  Every timing convention in the
literature leaves the time-zero argument untouched.  Even the strengthened win
condition ``occupy the closed neighbourhood $N[v]$'' survives: on $K_4$ that
set has $4$ vertices, still more than two cops can occupy; and the robber on a
cubic graph always retains a legal move to a cop-free neighbour, so a
``robber stuck'' loss condition also never triggers with two cops.  Finally,
$K_4$ is simple and cubic by any textbook definition, so no reading of
``cubic graph'' avoids the counterexample.
\end{remark}

\section{Formal verification}

The Lean~4 formalisation in \texttt{lean4/Main.lean} uses only the kernel and
core (no Mathlib, no \texttt{sorry}).  All definitions are pure structural
matches over explicit adjacency tables, and every statement is a Boolean
equation between closed computable expressions proved by \texttt{rfl}, i.e.\
by kernel reduction alone:

\begin{itemize}
\item \texttt{surrounded nbr cops v} --- the robber on $v$ is surrounded iff
  every neighbour of $v$ lies in \texttt{cops};
\item \texttt{twoCopTable nbr n} --- exhaustive over every robber vertex
  $v<n$ and every ordered pair of cop positions $(c_1,c_2)$, $c_1,c_2<n$
  (stacked cops included): the surround condition never holds;
\item \texttt{cubicTable nbr n} / \texttt{symmetricTable nbr n} --- sanity
  checks that the explicit tables define simple symmetric cubic graphs;
\item \texttt{K4\_two\_cops\_never\_surround}, \texttt{K4\_three\_cops\_surround}
  --- $\sigma(K_4)\ge 3$ and $\sigma(K_4)\le 3$, i.e.\ $\sigma(K_4)=3$;
\item \texttt{Petersen\_two\_cops\_never\_surround},
  \texttt{Q3\_two\_cops\_never\_surround},
  \texttt{K33\_two\_cops\_never\_surround},
  \texttt{prism\_two\_cops\_never\_surround} --- the Petersen graph, $Q_3$,
  $K_{3,3}$ and the triangular prism are likewise not 2-surroundable;
\item \texttt{all\_five\_two\_cops\_never\_surround} --- all five at once.
\end{itemize}

Running \texttt{lake env lean Check.lean} re-evaluates every audited quantity
(all come out \texttt{true}) and prints \texttt{\#print axioms} for all
seventeen theorems: \emph{every one of them reports ``does not depend on any
axioms''} --- no \texttt{propext}, no \texttt{Classical.choice}, no
\texttt{Quot.sound}, no \texttt{sorryAx}.  The script
\texttt{reproduce.py} (standard library only) repeats the sanity checks, the
exhaustive 2-cop tables, the $K_4$ witness and the retrograde game search, and
exits with status $0$ on success.

\begin{thebibliography}{9}
\bibitem{BCCDFJP}
A.~C.~Burgess, R.~A.~Cameron, N.~E.~Clarke, P.~Danziger, S.~Finbow,
C.~W.~Jones, D.~A.~Pike,
\emph{Cops that surround a robber},
Discrete Applied Mathematics (2021),
\url{https://doi.org/10.1016/j.dam.2020.06.019};
arXiv:\url{https://arxiv.org/abs/1910.14200}.
\end{thebibliography}

\end{document}
